Let $G_1,\ldots,G_m$ denote a finite collection of guardrails. For each
guardrail define a binary failure event
\[
  Z_i \in \{0,1\}, \qquad Z_i = 1 \text{ means guardrail failure or unsafe pass.}
\]
The unknown object is the atom distribution
\[
  \pi(z) = \Prob(Z=z), \qquad z \in \{0,1\}^m.
\]
The atom ordering used by the implementation is deterministic and little
endian, but the mathematical claims are invariant to relabeling.

Evaluation evidence may include singleton failure marginals
\[
  p_i = \Prob(Z_i=1)
\]
and optional pairwise overlap constraints
\[
  q_{ij} = \Prob(Z_i=1, Z_j=1),
\]
possibly as intervals. More generally, the implementation accepts linear
constraints of the form $a^\top \pi \le b$, $a^\top \pi \ge b$, or
$a^\top \pi=b$.

A composed safety question is a declared event $\phi:\{0,1\}^m\to\{0,1\}$.
Examples include at least one selected guardrail failing, all selected
guardrails failing, or a custom event over the finite atom space. Because
$\E_\pi[\phi(Z)]$ is linear in $\pi$, composition bounds are linear programs
over a finite simplex.

The distinction between mathematical and empirical assumptions matters. The LP
can verify what follows from supplied rates and constraints. It cannot verify
that those rates were measured on a representative dataset, that labels are
correct, or that future deployment prompts follow the evaluated population.
